\documentclass[11pt]{article}
% Shared preamble for all daily exercises
% Update this file to change settings (padding, parindent, etc.) for all exercises

\usepackage[margin=0.75in]{geometry}
\usepackage{amsmath}
\usepackage{amssymb}

% Custom settings
\setlength{\parindent}{0pt}
\setlength{\parskip}{0.2cm}

\title{Daily Exercise}
\author{Dhruman Gupta}
\date{24th March}

\begin{document}

% Common header for daily exercises
% This file can be included in other LaTeX documents

\maketitle

\noindent
\textbf{Course:} CS-3330, Theory of Computation

\vspace{0.2cm}

\noindent
\textbf{Question:}\noindent 
$L = \{\langle M,w,n\rangle \mid M \text{ halts on } w \text{ in at most } 2^n \text{ steps}\}$. Prove or disprove the decidability of membership in $L$.

\textbf{Answer:}

This language is decidable.

Let $D$ be a turing machine that takes input $\langle M,w,n\rangle$. It will do the following:

\begin{itemize}
    \item Compute $2^n$.
    \item Simulate $M$ on input $w$ for at most $2^n$ steps.
    \item If $M$ halts within those $2^n$ steps, accept.
    \item Otherwise, reject.
\end{itemize}

This is a valid decider because it always halts. The first step halts, and the simulation only runs for a finite number of steps, namely $2^n$.

Now, if $\langle M,w,n\rangle \in L$, then by definition $M$ halts on $w$ in at most $2^n$ steps. So the simulation will see this and accept.

If $\langle M,w,n\rangle \notin L$, then either $M$ does not halt on $w$, or it halts only after more than $2^n$ steps. In either case, the simulation will not see $M$ halt within $2^n$ steps, and so it rejects.

Hence, $D$ decides $L$. Therefore membership in $L$ is decidable.

\end{document}
